\section{Результаты}

\subsection*{Количество документов}

Объём корпуса оценён по счётчикам алфавитных категорий русской Википедии
(таблица~\ref{tab:cats}). Вложенные подкатегории не обходились, поэтому это нижняя оценка.

\begin{table}[h]
\centering
\begin{tabular}{|l|r|}
\hline
\textbf{Категория} & \textbf{Статей} \\
\hline
Фильмы по алфавиту & 38\,118 \\
Актёры по алфавиту & 15\,563 \\
Сценаристы по алфавиту & 6\,735 \\
Кинорежиссёры по алфавиту & 5\,502 \\
Телесериалы по алфавиту & 5\,148 \\
Мультфильмы по алфавиту & 4\,322 \\
\hline
\textbf{Итого} & \textbf{75\,388} \\
\hline
\end{tabular}
\caption{Объём кинотематики в русской Википедии}
\label{tab:cats}
\end{table}

Требование в 30\,000 документов перекрыто больше чем вдвое одним только первым
источником. Точное число документов на Кино-Театр.Ру сайт не публикует, но по номерам
страниц речь идёт как минимум о десятках тысяч записей.

\subsection*{Размеры документов}

Измерения по 12 сохранённым страницам приведены в таблице~\ref{tab:sizes}. Под
\enquote{сырым} размером понимается HTML-файл целиком, под текстом --- то, что осталось
после выделения.

\begin{table}[h]
\centering
\begin{tabular}{|l|r|r|r|r|}
\hline
\textbf{Документ} & \textbf{Сырой, КБ} & \textbf{Символов} & \textbf{Слов} & \textbf{Доля} \\
\hline
\multicolumn{5}{|l|}{\textit{ru.wikipedia.org}} \\
\hline
Во все тяжкие (сериал)      & 1152,3 & 160\,967 & 18\,132 & 18,9\,\% \\
Янковский, Олег Иванович    & 1011,4 & 134\,491 & 17\,757 & 21,5\,\% \\
Тарковский, Андрей Арс.     &  706,1 &  65\,359 &  9\,566 & 15,2\,\% \\
Криминальное чтиво          &  633,4 &  42\,945 &  6\,860 & 10,3\,\% \\
Сталкер (фильм)             &  523,3 &  45\,532 &  6\,965 & 14,6\,\% \\
13 дней во Франции          &  119,9 &     745 &     133 &  1,0\,\% \\
24 недели                   &  100,4 &     873 &     120 &  1,3\,\% \\
\hline
\multicolumn{5}{|l|}{\textit{www.kino-teatr.ru}} \\
\hline
Статья \enquote{Место встречи\ldots} & 105,3 & 18\,950 & 2\,711 & 32,0\,\% \\
Обзор сериала \enquote{Число зверя}  &  92,5 & 11\,433 & 1\,650 & 21,9\,\% \\
Биография: А. Мельникова             &  81,8 &  9\,374 & 1\,265 & 20,1\,\% \\
Карточка фильма                      &  82,1 &  8\,181 & 1\,144 & 17,4\,\% \\
Биография: Д. Мухамадеев             &  80,3 &  8\,320 & 1\,137 & 18,2\,\% \\
\hline
\end{tabular}
\caption{Размеры \enquote{сырых} документов и выделенного текста}
\label{tab:sizes}
\end{table}

Средние значения сведены в таблицу~\ref{tab:avg}.

\begin{table}[h]
\centering
\begin{tabular}{|l|r|r|}
\hline
\textbf{Показатель} & \textbf{Википедия} & \textbf{Кино-Театр.Ру} \\
\hline
Документов в выборке                          &       7 &       5 \\
Средний размер \enquote{сырого} документа, КБ &   606,7 &    88,4 \\
Средний объём текста, символов                & 64\,416 & 11\,252 \\
Средний объём текста, слов                    &  8\,505 &  1\,581 \\
Средняя доля полезного текста                 & 11,8\,\% & 21,9\,\% \\
\hline
\end{tabular}
\caption{Средние характеристики документа}
\label{tab:avg}
\end{table}

Самое заметное здесь --- доля полезного текста: у Википедии в среднем 12\,\%, а у
статей-заглушек всего 1\,\%. Причина в оформлении: тема встраивает в каждую страницу
около полумегабайта меню и скриптов независимо от длины статьи. Поэтому судить об объёме
будущего индекса по весу скачанных данных нельзя.

\pagebreak
